% ===========================================================================
%  07_discussion.tex -- 适用边界、测量不确定度、资源预算与展望
% ===========================================================================
\section{讨论}\label{sec:discuss}

第~\ref{sec:verify} 节把「算出来的量」与「它应该等于的值」逐点比对，
第~\ref{sec:results} 节把「跑起来时画面长什么样、代价是多少」变成可复现的测量。
前两节刻意把范围收窄到\emph{可验证}的那一部分；本节要补齐它的另一面：
哪些结论只在什么前提下成立，哪些近似会把误差藏到看不见的地方，
测量本身的不确定度有多大，以及在什么条件下本文的模型会失效。
我们把这一节写成一份「使用说明 + 失效清单」，而不是一份自我评价——
因为一个模拟器真正有用的部分，恰恰是它知道自己什么时候不能用的部分。

% --------------------------------------------------------------------------- %
\subsection{物理模型的适用边界：从施瓦西到 Kerr}\label{sec:discuss-phys}

\paragraph{本模拟器只解一个时空。}
全文所有的公式、验证与图像都建立在静态、不带电、球对称的
\emph{施瓦西}度规之上（式~\eqref{eq:schw}）。这不是一个技术选择，
而是一个物理假设：它要求中心天体角动量为零。真实的天体黑洞几乎
一定在转：对恒星质量黑洞，观测到的自旋参数 $a_\ast\equiv J/M^2$
可以在 $0.05$ 到 $0.98$ 之间取值~\cite{thorne1974}；
M87$^\ast$ 的 EHT 结果给出的自旋区间同样远离零~\cite{eht2019}。
因此本文渲染器的正确读法是「一个自旋为零的黑洞看起来是什么样」，
而不是「黑洞看起来是什么样」。

\paragraph{自旋会改变什么。}
把 $a_\ast$ 从 $0$ 提到 $0.9$（Kerr 度规、Boyer--Lindquist 坐标）
至少改变四类可观测的结构，且每一种都超出本文模型的能力：

\begin{enumerate}[leftmargin=1.8em,itemsep=2pt,topsep=3pt]
  \item \textbf{最内稳定圆轨道的位置。}施瓦西时空为
        $\bisco=6\M$；Kerr 时空正转时为
        $\bisco/\M=3+Z_2-\sqrt{(3-Z_1)(3+Z_1+2Z_2)}$，
        在 $a_\ast=0.9$ 处约为 $2.32\M$。盘的内边缘因此向里移动，
        而 Page--Thorne 通量正是以 $\bisco$ 为零扭矩边界积分的
        （式~\eqref{eq:Bclosed}）。在固定质量吸积率下，内区可发光
        面积的量级变化为 $(6/2.32)^2\approx6.7$ 倍，这已经足以
        改变整幅图像的亮度分布与峰值温度。
  \item \textbf{阴影的形状。}施瓦西阴影是严格的圆，角半径为
        $\sin\psi_{\rm sh}=\bc\sqrt{1-\rs/r_0}/r_0$
        （式~\eqref{eq:shangle}），本文正是依靠这个闭式把阴影半径
        验证到 $10^{-12}$ 量级。Kerr 阴影由 Bardeen~\cite{bardeen1973}
        给出的临界曲线界定，对 $a_\ast=0.9$ 而言它明显偏离圆，
        呈经典的「D 形」。EHT 的定量结论中有一项就是
        「阴影对圆的偏离不超过 $10\%$」~\cite{eht2019,eht2022}——
        这是一个\emph{只有带自旋模型才能讨论}的量，本文的渲染器
        连表达它的自由度都没有。
  \item \textbf{光子环的干涉及其 Lyapunov 指数。}近临界光线的
        绕行圈数在 Kerr 中也按对数发散，但发散系数由光子壳层
        （photon shell）的 Lyapunov 指数决定，而后者随自旋与
        轨道倾角变化~\cite{johnson2020,gralla2020}。
        第~\ref{sec:criticalband} 节测到的 $-2$ 与 $-1$ 两个斜率
        （式~\eqref{eq:slopekappa}--\eqref{eq:slopedx}）
        是施瓦西时空的普适常数，不能直接搬到 Kerr。
  \item \textbf{参考系拖曳。}能层（ergosphere）与
        Lense--Thirring 效应会在近场改变静态观测者标架的存在域：
        在 Kerr 中，$r<2\M$ 区域内不存在静态观测者，
        而本文的相机可以在任意 $r_0>\rs$ 处放置
        （第~\ref{sec:r0scan} 节的扫描一直推到 $r_0=8\M$，
        这已是采用静态标架能够靠近的极限附近）。
\end{enumerate}

\paragraph{为什么仍然选择施瓦西。}
理由不是计算便宜，而是\emph{可证伪}。施瓦西时空有四个本文实际用到的
解析靶子：临界冲击参数 $\bc=3\sqrt3\,\M$（式~\eqref{eq:bc}）、
光子球半径 $\bph=3\M$（式~\eqref{eq:photonsphere}）、
阴影角半径的闭式（式~\eqref{eq:shangle}）与弱场偏折角
$\alpha=4\M/b$（式~\eqref{eq:deflection}）。这四条使得
第~\ref{sec:verify} 节的六组实验都有 $10^{-6}$ 甚至 $10^{-12}$
量级的独立参照；换成 Kerr，能给出同等强度的闭式靶子只剩
Carter 常数与椭圆积分表示，验证链条会立刻变得依赖实现本身。
换言之，本文用「减少一个物理参数」换来了「验证误差预算小三个量级」。

\paragraph{通往 Kerr 的路径与代价。}
把模拟器扩展到 Kerr 在原理上是标准的~\cite{chandrasekhar1983,dexter2009}：
在 Boyer--Lindquist 坐标下运动仍然可分离，由 Carter 常数
$\mathcal{Q}$ 把径向与极角方向的方程解耦，每一条光线由
$(\lambda,\eta)$ 两个守恒量参数化。实现上有两个必须更换的部件：
（i）轨道平面的约化（第~\ref{sec:plane} 节）不再成立，
必须同时积分 $(r,\theta)$ 两个方向的运动并处理两者的转向点，
状态量由二维 $(u,u')$ 变成四维；
（ii）终止与步长判据必须同时覆盖径向与极角的转折
（$R(r)=0$、$\Theta(\theta)=0$ 的零点搜索）。
两者叠加的代价可以直接估计：每次 RK4 求值的工作量约翻倍，
近转向点处所需的步数增加，因此在同一像素数与同一容差下，
整条管线的成本大致上升到 $2$--$4$ 倍。这个倍数是在
本文的代价模型（式~\eqref{eq:costpix}）之上直接乘上去的，
不需要新的拟合——这也正是把代价写成「每像素成本 $\times$ 像素数」
这一形式的用处。

% --------------------------------------------------------------------------- %
\subsection{吸积盘与辐射转移的近似}\label{sec:discuss-rt}

第~\ref{sec:disk} 节给出的盘模型是一个\emph{几何薄、光学厚、
无自吸收稳态盘}，其通量剖面来自 Page--Thorne~\cite{pt1974}
的时间平均扭矩积分（式~\eqref{eq:Bclosed}）。把这个模型当作
「真实吸积盘的照片」会产生系统偏差，偏差来源可以逐条列出。

\paragraph{（一）薄盘假设。}
盘被放在赤道面 $z=0$ 上，不存在垂向的结构解。真实薄盘的
垂向厚度 $H/r$ 约为 $10^{-3}$--$10^{-1}$，且随半径增长
（$H\propto r^{9/8}$）。本文在体积辐射转移中人为加了一个
垂向高斯权重（式~\eqref{eq:gauss}）与湍流调制
（式~\eqref{eq:turbmod}），它让近侧盘的边缘不再是数学上的
一条线，但这只是\emph{视觉上}的软化，不是流体静力学平衡解。
与发射面的垂向位置直接相关的可观测后果——例如边缘被挡的
临界倾角、以及吸积盘对自身辐射的遮挡——因此只被近似描述。

\paragraph{（二）通量剖面的模型依赖。}
Page--Thorne 通量的推导依赖三个前提：稳态、无粘性扭矩
在 $\bisco$ 处为零、以及角动量输运可以写成局部的一维积分。
真实盘中的磁转动不稳定性、盘风与喷流都会破坏其中至少一条。
本文把 PT 与 Shakura--Sunyaev~\cite{ss1973} 剖面做成
运行时开关（式~\eqref{eq:ss} 与式~\eqref{eq:pt}），
正是为了把「模型不确定性」暴露成一个可见的自由度：
两套剖面在峰值区相差 $3$ 倍以上（图~\ref{fig:07}），
这个差距的量级给出了通量模型的理论不确定度下限。
值得强调的是，这个 $3$ 倍差距\emph{不是}数值误差，
而是模型选择的结果；任何引用本文图像的定量结论都必须
同时说明使用的是哪一套剖面。

\paragraph{（三）黑体配色与光谱硬化。}
本文由 $T_{\rm eff}(r)$ 经 Planck 谱与 CIE 1931 色度轨迹
映射到显示颜色（图~\ref{fig:10}）。真实盘的出射谱不是
该温度的黑体：电子散射会使其在约化频率上发生硬化，
通常用一个约 $1.7$--$2.0$ 的谱硬化因子 $f_{\rm col}$
来参数化。本文没有引入 $f_{\rm col}$。其后果是：
在固定光度下，可见光通道拿到的相对亮度与色相会有
系统偏移，偏移量级由 $f_{\rm col}^4$ 决定，即
$8$--$16$ 倍以内的亮度重分配。这就是第~\ref{sec:sweep}
节中「温度扫描不改变相对论增亮图案、只改变色标」这一
结论之所以安全的原因：几何结构由 $g$ 因子决定，
与 $f_{\rm col}$ 无关；而绝对亮度则不安全。

\paragraph{（四）前向体积转移的边界。}
第~\ref{sec:rtnum} 节实现的是一维前向（front-to-back）
累积（式~\eqref{eq:front2back}），近盘处用 $K=5$ 个子段
细化，光学深度用 $\tau=\int\kappa\rho\,\dd l$ 的
分段近似（式~\eqref{eq:avgw}）。它\emph{不是}辐射转移方程的
自洽解，因为它缺三样东西：

\begin{itemize}[leftmargin=1.6em,itemsep=1pt,topsep=3pt]
  \item \textbf{自吸收与再发射。}累积式只把发射项
        $S\,(1-e^{-\dd\tau})$ 加权叠加，不求解
        $\dd I/\dd\tau=-I+S$ 的完整边值问题，也不让
        介质吸收自己发出的光后再辐射。对光学厚的盘内区，
        这等价于假设源函数在局部等于 Planck 函数，
        恰好是薄盘模型的常用近似，因此误差可控；
        对光学薄的外区与盘冕，它会把亮度高估，
        因为真实的薄介质是「发射减吸收」而不是纯发射。
  \item \textbf{散射与偏振。}电子散射被归入灰度消光系数，
        不作角向重分布，也不产生偏振。偏振是 EHT 的
        核心观测量之一~\cite{eht2022,grtrans2016}，
        本文完全没有涉及。
  \item \textbf{康普顿化。}高温冕中的光子会通过逆康普顿
        散射获得能量，产生幂律硬尾~\cite{koral2013}。
        本文的 X 射线预设（图~\ref{fig:19} 第 5 幅）
        只是把峰值温度推到 $10^{7}$\,K 后观察
        Rayleigh--Jeans 尾部落在可见光通道的结果，
        \emph{不}包含任何非热成分。
\end{itemize}

\paragraph{（五）$g^{4}$ 集束律的应用对象。}
这里有一处容易写错、也值得在本文中明确的地方：频率积分
强度的洛伦兹不变量是 $I_\nu/\nu^3$，由此得到的
频率积分（bolometric）强度变换是
$I_{\rm bol}^{\rm obs}=\gfac^{4}I_{\rm bol}^{\rm em}$
（式~\eqref{eq:I4}），而单色强度是
$I_\nu^{\rm obs}=\gfac^{3}I_\nu^{\rm em}$。
本文的着色器给出的是\emph{bolometric} 量：先由 PT 通量得到
局部有效温度，把 $T^4\propto F$ 作为发射率，再整体乘 $\gfac^{4}$
（第~\ref{sec:impl} 节所列的 \texttt{diskSource()} 实现），
颜色温度则按 $T_{\rm obs}=\gfac\,T_{\rm em}$
（式~\eqref{eq:Tobs}）偏移。两者是自洽的：温度按 $\gfac$
缩放、bolometric 强度按 $\gfac^4$ 缩放，正好等于
$T^4$ 关系。若把 $\gfac^3$ 误用到 bolometric 量上，
会得到偏暗的像；若把 $\gfac^4$ 误用到单色谱上，
会得到偏亮的像。第~\ref{sec:shift} 节的推导与
第~\ref{sec:impl} 节的实现都以这一条为准。

\paragraph{（六）时间平均。}
PT 通量是对粘性耗散做时间平均之后的稳态解，因此本文的图像
代表「一个稳态盘的期望像」，而不是任何一个瞬时快照。
真实盘的湍流会把亮度在时标 $r/c_s$ 上波动，GRMHD 快照中的
瞬时结构（热斑、磁通量管）可以显著改变局部亮度~\cite{koral2013,ipole2016}。
本文的湍流调制（式~\eqref{eq:turbmod}）是程序化噪声，
它给出空间上的非均匀感，但不携带任何流体动力学内容。

% --------------------------------------------------------------------------- %
\subsection{背景星场：结构可验证，光度不可验证}\label{sec:discuss-stars}

第~\ref{sec:stars} 节的星场由程序化噪声与幂律星等分布生成
（式~\eqref{eq:starmag}），不对应任何真实星表。这一点的
后果必须区分清楚：

\begin{itemize}[leftmargin=1.6em,itemsep=1pt,topsep=3pt]
  \item \textbf{可验证的部分是几何。}星场的像在引力透镜下
        被搬到哪、被拉伸多少、二级像在哪里出现，全部由
        光线的出射方向决定，而出射方向在
        第~\ref{sec:v5} 节被验证到 $0.021$\,px 以内。
        图~\ref{fig:02} 的偏折曲线与图~\ref{fig:15} 的
        「关闭引力 / 开启引力」对照说明：透镜结构本身是
        正确几何的产物。
  \item \textbf{不可验证的部分是光度。}任何一个具体像素的
        亮度取决于那颗（虚构的）星的星等、能谱与距离，
        而这些量没有任何独立来源。因此本文不使用星场做
        任何光度学结论；爱因斯坦环的\emph{位置}可以引用，
        环内外的\emph{亮度比}不可以。
\end{itemize}

一个由此产生的实用建议是：如果要把该渲染器用于教学演示
「引力透镜」，应当使用关闭吸积盘的预设（\emph{纯引力透镜}，
图~\ref{fig:19} 第 4 幅），此时画面上除阴影外唯一的
物理内容是星场的畸变，视觉与几何一一对应，不会与盘的
辐射模型混淆。

% --------------------------------------------------------------------------- %
\subsection{计时量化与测量不确定度}\label{sec:discuss-quant}

第~\ref{sec:perf} 节的代价模型是本文唯一一个\emph{非物理}的
定量结论，因此它的不确定度需要单独交代。我们在这里用
原始基准记录（\texttt{render\_benchmark.json}，36 组配置、
每组 $n=7$ 次采样）逐项量化。

\paragraph{同一配置内部的离散远大于配置之间的差异。}
对 36 组配置分别计算 $t_{\max}/t_{\min}$：最小值 $3.26$、
中位数 $5.55$、最大值 $24.14$；其中 $27/36$ 组的比值超过 $5$。
作为最极端的一例，$400\times240$、$75$ 步的那一组
$t_{\min}=62.1$\,ms 而 $t_{\max}=1189$\,ms，相差 $19$ 倍。
这个离散\emph{不}来自测地线积分的随机性——积分器是确定性的，
同一配置的输出逐位相同（第~\ref{sec:v6} 节已验证）——
所以它只能来自计时通道本身。

\paragraph{四个可分辨的来源。}

\begin{enumerate}[leftmargin=1.8em,itemsep=2pt,topsep=3pt]
  \item \textbf{GPU 队列深度与返回时序。}测量的是相邻两次
        \texttt{requestAnimationFrame} 回调的时间间隔。着色器
        提交后回调即返回，若主机侧有反压（back-pressure），
        队列深度会把一次长时间停顿折算进某一帧，表现为
        一个孤立的大 $t_{\max}$，而中位数几乎不动——
        实测正是如此（$t_{\max}/t_{\rm median}$ 的最大值
        $23.77$，与 $t_{\max}/t_{\min}$ 同阶）。
  \item \textbf{热漂移与共享功耗预算。}被测机器是一台只有
        集成显卡的机器，CPU 与核显共享同一功耗墙。一次完整
        扫描持续数分钟，期间封装温度上升、时钟回落，
        于是同一配置在扫描早期与晚期的耗时不同。
        这解释了 $t_{\rm median}$ 也偏离 $t_{\min}$
        数倍的现象。
  \item \textbf{宿主进程的其它工作。}扫描脚本在同一进程外
        同时在做屏幕捕获与 JSON 写入，任何一次数毫秒到
        数百毫秒的系统级抢占都会落在某一帧上。
        $t_{\max}\approx1.2$--$1.5$\,s 的量级与
        一次捕获/写盘停顿相符。
  \item \textbf{计时器量化。}WebView2 出于安全考虑对
        \texttt{performance.now()} 有钳制，粒度在
        $100\,\mu$s 到 $1$\,ms。本文最短的帧是 $62$\,ms，
        所以这一步在\emph{本文的测量范围}内不是主导项，
        但它设定了「无法测量到优于 $\sim1$\,ms」的硬下限；
        对更小的视口或更少的步数（帧时降到毫秒量级），
        它将成为主要误差源。
\end{enumerate}

\paragraph{因此我们报告的是下界。}
统计量选择上，我们用「最干净子集的最小值」而非均值或中位数
来估计固有成本：把 $\{t_i\}$ 排序后取 $t_{\min}$，
再对 12 个像素律测点做一次最小二乘并剔除一个最坏点。
$t_{\min}$ 是「没有受到停顿污染的一帧」，因而它系统性
\emph{偏小}。这意味着式~\eqref{eq:costpix} 给出的
$k=6.298\times10^{-4}$\,ms/px 是真实每像素成本的\emph{下界}，
而代价模型在「低端硬件上够不够用」这一类问题上是保守的：
真实成本只会更高，不会更低。

\paragraph{1.9\% 的口径差异说明了精度上限。}
同一点（$1000\times600$、$tol=0.03$）在
表~\ref{tab:perfsteps} 与表~\ref{tab:perftol} 中分别给出
$388.3$\,ms 与 $395.8$\,ms，差 $1.9\%$；而模型的残差
均方根是 $2.37\%$。两者同阶说明：本文的代价模型
只有效到约 $2\%$——任何小于 $2\%$ 的模型差异都不应被
解释为物理效应。$p=0.912$ 与 $1$ 的偏离（$8.8\%$）
则明确大于这一噪声水平，因此「$p$ 显著小于 $1$」是
一个成立、且如第~\ref{sec:perf} 节所述有物理原因
（多数像素提前终止）的结论。

\paragraph{元数据粒度这一处必须诚实。}
基准记录里的容差字段只取过 $0.06/0.03/0.01$ 三个值。
表~\ref{tab:perftol} 中 $tol=0.0150$ 与 $tol=0.0075$
两档的行，来自扫描脚本的\emph{档位定义}，而不是逐条
写进记录文件的物理设置。其中 $tol=0.0075$ 的四行被
$n_{\max}=4096$ 截断，本来就不进入拟合；$tol=0.0150$
的四行耗时是真实测量，但它们与 $tol=0.01$ 同属
「预算够用」区，我们只声称它们属于该区，不声称它们
各自对应一个独立的、被逐条记录的 $0.0150$ 设置。
把这条写进论文的原因是：如果不写，一个读者按
表~\ref{tab:perftol} 去复现 $tol=0.0150$ 的四个耗时，
可能得到与其名义档位不符的结果，而这会被误读为
本文的数值不可复现。实际的不可复现来源只是元数据。

\paragraph{几何量的量化台阶，以及一个被“定量幻觉”掩盖的缺陷。}
阴影半径探针的分层覆盖率量化带固定为
$\sigma_{\rm px}=1/(2\sqrt2\cdot24)=0.01473$\,px
（式~\eqref{eq:subpix}），与视口和 $r_0$ 都无关；
被测的像素半径却按 $1/r_0$ 缩小。因此人们很容易预期：
相对误差/信号比会随 $r_0$ 增长，小信号端测量最差。
初版探针的扫描\emph{看起来}正是这样——残差虽小，但当
$r_0$ 增大时读数逐渐靠拢整数格点，被解释为「量化误差开始主导」。
这个解释听上去完全合理，却掩盖了一个真实缺陷：
估计器在边缘落到外侧像素时被钳位，退化成整数扫描
（式~\eqref{eq:subpix} 的讨论）。

把它纠正之后，观测到的行为与上述预期\emph{不一样}：
第~\ref{sec:r0scan} 节的七个档位跨越 $r_0/M\in[8,150]$、
像素半径 $552\to28$\,px，残差却稳定停在百分之几个像素，
既不随 $r_0$ 放大，也不随信号缩小而相对恶化（表~\ref{tab:r0scan}）。
正确的读法是把误差归因于\emph{读数格点}而不是被测物理量：
量化带 $\sigma_{\rm px}$ 与掩膜的系统性亚像素项相加，
给出一条与 $r_0$ 近似无关的误差地板。这条经验值得写下来，
因为它解释了为什么本文把探针精度写成「限制项清单」而不是
一个随 $r_0$ 变化的误差模型。

更一般的教训在方法论层面：一个\emph{小}残差并不蕴含
一个\emph{正确}的估计器。单点测量无法区分
「物理模型偏了百分之一」与「估计器悄悄退化成整数扫描」，
因为两者在同一视口下给出同量级的残差；只有多变视口、
多半径的残差\emph{形状}才能把二者分开。
这也是本文所有定量验证都坚持跨参数扫描、
而不是单点报数的原因。

\paragraph{8\,bit 量化曾经把累积管线整体废掉。}
第~\ref{sec:v6} 节的核心发现是：在修正之前，抖动相位与
累积历元不同步，于是渐进累积在 $8$\,bit 量化台阶上饱和，
亮度地板约为 $4\times10^{-3}$，继续累积不再改善图像。
修正后残差按 $n^{-0.29}$ 下降。指数 $0.29$ 与理想蒙特卡洛的
$n^{-0.5}$ 有差距，这个差距本身是有意义的：它说明误差中
有一部分是\emph{系统性的}，不会随采样数下降——
具体就是 RK4 的截断误差、轨道平面约化的近似与
覆盖率量化（式~\eqref{eq:subpix}）这三项，
它们对任何有限 $n$ 都不被平均掉。因此渐进累积的
正确使用方式是「把随机采样误差压到系统性误差之下即停止」，
而不是「无限累积」。

\paragraph{小结：哪些数字可以引用。}
综合以上，本文的定量结论按可信度分为三层：
（i）\textbf{可无条件引用}——所有与解析闭式比对的量
（$\bc$、阴影角、偏折角、步长收敛、逃逸半径误差），
它们的残差在 $10^{-6}$--$10^{-14}$，远离任何测量噪声；
（ii）\textbf{可引用但需注明口径}——代价模型的 $k$、$c$、$p$
与 $r_0$ 扫描偏差，精度约 $2\%$，且是低端硬件的下界；
（iii）\textbf{仅作定性}——所有与绝对亮度、光谱形状、
湍流结构有关的描述，因为它们依赖盘模型与配色近似
（第~\ref{sec:discuss-rt} 节）。

% --------------------------------------------------------------------------- %
\subsection{代价模型的适用边界与资源预算}\label{sec:discuss-perf}

\paragraph{线性律是局部展开，不是全局定律。}
式~\eqref{eq:costpix} 的拟合域是 $9.6\times10^4$ 到
$6.0\times10^5$ 个受追踪像素（$40$ 倍动态范围），
固定盘开启、固定后处理参数。把它外推到任意配置会遇到两处
结构性偏差：

\begin{itemize}[leftmargin=1.6em,itemsep=1pt,topsep=3pt]
  \item \textbf{固定开销 $c$ 依赖画布尺寸而非追踪像素数。}
        辉光金字塔与 ACES 是在\emph{画布}分辨率上运行的
        （第~\ref{sec:pipeline} 节），而在像素律扫描中
        画布尺寸恒定、只有离屏追踪目标在缩放，所以 $c$ 在
        那 12 个点上确实是常数。一旦用户改变画布尺寸，
        $c$ 会随之变化，线性形式应改写为
        $T\propto k_{\rm tr}N_{\rm tr}+k_{\rm post}N_{\rm canvas}+c_0$。
  \item \textbf{每像素成本 $k$ 依赖场景。}盘关闭时，
        近盘像素的求交与体积转移（第~\ref{sec:rtnum} 节）
        整段消失，$k$ 会显著下降；星场采样只有一次纹理
        查找（式~\eqref{eq:escapedir2}）。因此
        $k=6.298\times10^{-4}$\,ms/px 是\emph{盘开启}的
        $k$；对「纯引力透镜」预设应当重新标定。
\end{itemize}

\paragraph{步数预算与容差的耦合。}
第~\ref{sec:perf} 节的结论之一是 $n_{\max}\geq300$ 后
耗时不再有系统趋势（25 组比值在 $[0.895,1.089]$ 内）。
但这一饱和是在 $tol=0.03$ 附近测到的；把 $tol$ 降一个
量级，步数需求随之上升，$n_{\max}$ 的「够用」阈值也会
上升。表~\ref{tab:perftol} 中 $tol=0.0075$ 的四组撞上
$n_{\max}=4096$ 正是这一耦合的直接证据：在极小容差下，
预算本身成为主导误差，继续降容差只会增加未完成的工作。
因此默认值的选择（$n_{\max}=300$、$tol=0.02$--$0.03$）
是一个\emph{交互性优先}的决定，而「高分辨率静帧」预设
（$800$ 步、$tol=0.015$、\texttt{accumulate}）则是
\emph{正确性优先}的对照。

\paragraph{内存预算。}
渲染目标共 10 个：追踪目标、两份累积乒乓历史与三张半分辨率
辉光中间图、两张四分之一与两张八分之一分辨率的金字塔层
（第~\ref{sec:pipeline} 节，\texttt{allocateTargets()}）。
以追踪像素数 $A$ 为单位，总像素当量为
$3A+3(A/4)+2(A/16)+2(A/64)\approx3.91A$。
在 RGBA16F（$8$ 字节/像素）下，主机侧把
$A$ 限制在 $2.2\times10^6$ 以内（\texttt{MAX\_PIXELS}），
对应约 $69$\,MB 的显卡内存，加上默认帧缓冲后接近 $80$\,MB。
这个上限是刻意设的：更强的适配器上把 $A$ 放大到
$3$--$4$ 倍（例如 $4$K 画布）会让显存占用进入数百 MB 级，
在集成显卡上容易触发设备丢失而不是仅仅变慢。
\texttt{renderScale} 机制因此既是性能旋钮，也是稳定性旋钮：
它改的是被追踪的像素数 $A$，而后处理仍在画布分辨率上运行。

\paragraph{自适应的行为。}
\texttt{autoQuality} 在帧率低于 $24$\,fps 时按 $0.1$ 的
步长下调 \texttt{renderScale}（下限 $0.45$），高于 $58$\,fps
时按 $0.05$ 上调（上限 $1.0$），并用
\texttt{ResizeObserver} 与 \texttt{bh-viewport} 事件在
侧栏开合时重新分配目标（第~\ref{sec:ui} 节）。
这带来一个测量上的告诫：\emph{基准测试必须关闭自适应}，
否则扫描过程中的 \texttt{renderScale} 变化会把像素律
的两端混在一起。本文的所有性能数据都是在
\texttt{autoQuality} 关闭、\texttt{renderScale} 固定的
条件下采集的。

% --------------------------------------------------------------------------- %
\subsection{「视觉说服力」掩盖的四类缺陷}\label{sec:discuss-lessons}

写作本文的过程中，我们在这个模拟器里发现并修正了六处实现缺陷
（第~\ref{sec:intro} 节列出的清单）。它们的共同特征值得单独
提炼成一条方法论：\textbf{这六处缺陷没有一处能通过「看图」发现}。
每一张截图在缺陷存在时都同样「好看」，物理上却已经是错的。
按被掩盖的机制分类，它们是：

\begin{enumerate}[leftmargin=1.8em,itemsep=2pt,topsep=3pt]
  \item \textbf{状态残留型。}预设系统未做全量基线
        （第~\ref{sec:presets} 节）导致切换预设后残留上一组的
        参数：画面只是「略有不同」，而读者无法判断那点不同
        来自参数还是来自残留。修正方式是让每个预设都从
        一份完整的基线展开，而不是只覆盖差异字段。
  \item \textbf{键空间型。}通量模型开关 \texttt{fluxModel}
        没有进入累积键，于是 PT 与 SS 两套剖面在渐进累积中
        被混叠成同一张图：两套剖面各自都正确，混合的结果
        没有对应的物理对象。修正方式是把「影响像素值的
        每一个自由度」都纳入历史重建的判据。
  \item \textbf{语义漂移型。}盘温度参数 \texttt{diskTemp}
        在显示模式下被当作色温使用，而物理模式下它是
        有效温度：同一个数字在两处代表不同的量，界面却
        只有一个滑块。这类缺陷在单一模式下永远测不出来，
        只有把两种模式并排比较才暴露。
  \item \textbf{时序假设型。}抖动种子从未随时间推进，
        且 $8$\,bit 抖动相位与累积历元不同步
        （第~\ref{sec:jitter} 节）。后果是第~\ref{sec:v6} 节
        量到的量化地板：累积看似在工作（图像确实在变化），
        实际在几步之后就停在量化台阶上不再收敛。
\end{enumerate}

这条方法论的实践含义是：对一个声称「物理正确」的渲染器，
仅仅展示图像是不够的；必须给出\emph{可证伪的对照量}，
并让每一个影响输出的自由度都进入可复现键。
本文第~\ref{sec:verify} 节的六组实验与
第~\ref{sec:results} 节的三段模型，都是这条要求的产物。
反过来说，本项目最诚实的自我评价是：它之所以能被引用，
不是因为画面好看，而是因为它报告了自己曾经的错误、
以及这些错误被发现的判据。

% --------------------------------------------------------------------------- %
\subsection{与既有工具的相对定位}\label{sec:discuss-related}

按第~\ref{sec:related} 节的三层分类，本文位于「实时可视化层」，
但它同时借用了「解析/半解析层」的闭式靶子，并补齐了
该层通常缺失的误差预算。逐项对照如下。

\paragraph{相对 GPU 光线追踪器（GRay~\cite{gray2013}、
GRay2~\cite{gray2}、RAPTOR~\cite{raptor2019}）。}
这些实现的目标是 Kerr 时空上的高吞吐成像，功能与物理覆盖面
都超过本文：它们能处理自旋、能与 GRMHD 快照耦合、能在
$10^7$--$10^8$ 像素量级上批量出图。本文的方向不同：
用单一时空换取一条把误差压到 $10^{-12}$ 的验证链，
并把「代价随参数如何变化」写成闭式。
两者是互补的：本文的验证方法（闭式靶子 + 独立参考实现 +
亚像素反解）可以直接搬到那些实现在施瓦西极限下的
回归测试中。

\paragraph{相对半解析方法。}
这一类方法用分离变量加椭圆积分，在大偏折区精度极高，且在近临界处
仍然稳定；本文的 RK4 自适应积分在同样的区域需要
$10^3$ 步量级并把守恒量误差压到 $2.6\times10^{-6}$
（第~\ref{sec:v2} 节）。差别在于：半解析方法给出的是
单条光线的解，把它接到「每像素一次、要在 $16$\,ms 内
完成 $6\times10^5$ 次」的管线里需要一个适合 GPU 的
求值策略；本文选择的是可直接向量化的数值积分。
其中的代表实现是 geokerr~\cite{dexter2009} 与
Cunningham--Bardeen~\cite{cunningham1973} 的分离变量解。

\paragraph{相对 GRMHD 耦合管线（KORAL~\cite{koral2013}、
IPOLE~\cite{ipole2016}、grtrans~\cite{grtrans2016}）。}
它们的等离子体模型由第一性原理模拟给出，包含自吸收、
同步辐射与偏振；本文的盘是一个解析模型，
辐射转移是前向累积的近似（第~\ref{sec:discuss-rt} 节）。
这意味着：本文的图不能用于拟合观测数据，
而上述管线可以。本文能提供的是「在给定解析盘模型下，
几何与频移部分可以被验证到多准」这一问题的答案。

\paragraph{相对实时图形学实现（Müller 等~\cite{muller2010,muller2014}、
Ament 等~\cite{ament2015}）。}
Ament 等人系统分析了 GPU 上零测地线积分的稳定性与
性能边界，是本文第~\ref{sec:perf} 节最直接的参照。
本文与它的差别在于把性能分析进一步写成对用户的
参数（像素数、步数上限、容差）的显式模型，
并给出这一模型在低端硬件上的残差；
同时也把交互层的两个「本应平凡、实际极易出错」的
不变量（中心不变量、累积的逐位可复现性）写成了
可测的判据（第~\ref{sec:ui} 节、第~\ref{sec:v6} 节）。

\paragraph{一张定位表（用文字给出）。}
把上面的对照压成三条判据：本文的\textbf{强度}是
（i）施瓦西几何与阴影的闭式验证到 $10^{-12}$、
（ii）可复现的代价模型与逐位可复现的累积、
（iii）零依赖的便携式浏览器实现；
本文的\textbf{边界}是（i）无自旋、（ii）无自吸收/
偏振/康普顿、（iii）单一时空的解析盘模型；
本文\textbf{不建议}被用于的场合是（i）与 EHT 观测
做定量拟合、（ii）偏振或谱线预测、
（iii）任意自旋下的强场结构研究。

% --------------------------------------------------------------------------- %
\subsection{展望}\label{sec:discuss-future}

沿第~\ref{sec:discuss-phys}--\ref{sec:discuss-perf} 节的边界，
把本工作往下推的最短路径有四条，按「物理收益/实现代价」排序如下。

\paragraph{（一）Kerr 时空与 Carter 常数。}
这是收益最大的一步，因为它同时解除自旋、ISCO 位置与
阴影非圆性三个限制。实现路径见第~\ref{sec:discuss-phys} 节：
四维状态、$(r,\theta)$ 双向转向点、椭圆积分或 RK 的
混合求值。验证策略可以完全沿用本文的框架——
把 $\bc=3\sqrt3\M$ 的闭式换成 $a_\ast=0$ 时的极限值，
其余的验证实验（守恒量、亚像素阴影、逃逸方向、
累积收敛）在 $a_\ast\to0$ 时都应当退回本文测到的数字。
这是一条很强的回归判据：任何 Kerr 实现都应当通过
「$a_\ast=0$ 时复现本文全部表格」的测试。

\paragraph{（二）自吸收与康普顿化。}
把前向累积升级为求解 $\dd I/\dd\tau=-I+S$ 的自洽形式，
需要（i）一个频率网格或灰体近似下的谱硬化因子、
（ii）发射率与吸收系数的显式模型（自由—自由、
同步辐射、逆康普顿）、（iii）对光学薄区的正确边界条件。
收益是让 X 射线预设从「色标演示」变成可讨论的能谱。
代价是每次光线—盘交互的成本上升一个量级，
实时性需要用降分辨率 + 累积来换。

\paragraph{（三）WebGPU 计算着色器。}
测地线积分是天然并行的，且近临界曲线附近的像素需要
远多于平均的步数（第~\ref{sec:criticalband} 节
的对数发散正说明这一点）。在 WebGPU 上可以做两级策略：
第一遍用低步数预算判定哪些像素是「近临界」，
第二遍用计算着色器对这些像素做细分与高预算积分，
并把结果写回同一张浮点纹理。这一方向不改变物理，
但可能带来 $5$--$20$ 倍的有效吞吐，
使 Kerr + 自吸收在同样硬件上重新变得实时。
保留 WebGL2/WebView2 路径的价值在于可移植性：
任何 Windows 机器上的便携包都能直接双击运行，
不依赖用户的浏览器版本或额外的运行时。

\paragraph{（四）可微渲染与硬件重标定。}
如果给每一条光线的积分写出手写的伴随（adjoint）
递推，整条管线就可以对参数（自旋、倾角、内缘半径、
温度剖面）求梯度，从而把「渲染」变成「拟合」：
用观测图像的反演来标定盘参数。这在数学上是最吸引人的
方向，但它要求整套模型先具备第~\ref{sec:discuss-phys}
与第~\ref{sec:discuss-rt} 节所缺的自由度。与之并行的
是一项成本极低、收益明确的工作：把第~\ref{sec:perf} 节的
扫描在一张英伟达显卡上重跑一遍，用实测值替换
$k$、$c$ 与 $p$，并把结果与本文的
下界性质（第~\ref{sec:discuss-quant} 节）并列报告。
这样代价模型的适用范围会从「一张集成显卡」扩展到
「一类硬件」，而推导它的方法不变。

\paragraph{（五）可交付性本身。}
本文的模拟器被设计成\emph{可交付}的：便携文件夹内的
可执行文件双击即运行，不需要安装 Python、Node 或显卡驱动，
并且默认走 GPU 加速路径（第~\ref{sec:pack} 节）。
这一属性让它的验证结果可以被任何审稿人独立复现，
而不依赖某个集群或环境。我们认为，对一个教学与
可复现性导向的模拟器而言，这与其物理内容同等重要——
第~\ref{app:repro} 节给出的复现命令全部只依赖
仓库内的文件与一个 LaTeX 引擎。

\FloatBarrier
